% Generated draft: requires mathematical and source review.
\documentclass[11pt]{article}
\usepackage[margin=25mm]{geometry}
\usepackage{fontspec}
\setmainfont{Latin Modern Roman}
\usepackage{amsmath,amssymb,mathtools,unicode-math}
\setmathfont{Latin Modern Math}
\usepackage{enumitem,needspace,xurl,hyperref}
\setlength{\parindent}{0pt}
\setlength{\parskip}{5pt}
\setlength{\emergencystretch}{3em}
\newcommand{\N}{\mathbb{N}}
\newcommand{\Z}{\mathbb{Z}}
\newcommand{\Q}{\mathbb{Q}}
\newcommand{\R}{\mathbb{R}}
\newcommand{\C}{\mathbb{C}}
\newcommand{\F}{\mathbb{F}}
\newcommand{\A}{\mathbb{A}}
\newcommand{\PP}{\mathbb{P}}
\newcommand{\E}{\mathbb{E}}
\newcommand{\eps}{\varepsilon}
\newcommand{\dd}{\,\mathrm{d}}
\begin{document}
\section*{UnsolvedMath ID: 700224}
\textbf{Source problem code:} catalogue entry
\subsection{Definitions and mathematical statement}
For a prime number \(p\), a finite \(p\)-group is a finite group \(G\) with \(|G|=p^n\) for some integer \(n\geq 0\). Write \(\operatorname{Aut}(G)\) for its automorphism group and
\[
\operatorname{Inn}(G)=\{c_g:g\in G\},\qquad c_g(x)=gxg^{-1}\quad(x\in G),
\]
for its subgroup of inner automorphisms. An automorphism \(\alpha\) has order exactly \(p\) if \(\alpha^p=\operatorname{id}_G\) and \(\alpha\ne\operatorname{id}_G\).

The proposed assertion is: for every prime \(p\) and every finite non-abelian \(p\)-group \(G\), there exists \(\alpha\in\operatorname{Aut}(G)\) such that
\[
\alpha^p=\operatorname{id}_G,\qquad
\alpha\ne\operatorname{id}_G,\qquad
\alpha\notin\operatorname{Inn}(G).
\]
The supplied registry labels this problem open; that status is supplied source metadata and has not been independently verified.
\subsection{Short English statement}
Every finite non-abelian \(p\)-group is conjectured to have an automorphism of order exactly \(p\) that is not conjugation by any element of the group, for every prime \(p\).
\subsection{Sources}
\begin{itemize}
\item Local UnsolvedMath registry, supplied complete record: permanent ID \(700224\), exact source code \texttt{catalogue entry}, title \emph{Kourovka Notebook Problem 4.13}. The record labels the problem open and gives its update timestamp as \texttt{2026-10-16T14:00:00Z}; this timestamp is registry metadata, not an independently verified literature assessment.
\item \emph{The Kourovka Notebook (Unsolved Problems in Group Theory)}, Problem 4.13, fourth issue (1973); the supplied record identifies the twenty-first issue edition (2026), arXiv identifier \texttt{1401.0300v48}.
\item Proposed by Ya. G. Berkovich in 1973. The supplied collection description attributes the editing of Issues 1--21 (1965--2026) to E. I. Khukhro and V. D. Mazurov.
\end{itemize}
\end{document}
